\documentclass[a4paper,10pt,twoside,openany]{article}

\usepackage[lang=hebrew]{maths}
\usepackage{hebrewdoc}
\usepackage{stylish}
\usepackage{lipsum}
\let\bs\blacksquare

\setlength{\parindent}{0pt}

%%%%%%%%%%%%
% Styling %
%%%%%%%%%%%%

\usepackage{enumitem}

%%%%%%%%%%%%%
% Counters  %
%%%%%%%%%%%%%

\setcounter{section}{1}     
            
%BIBLIOGRAPHY
\usepackage[
backend=biber,
style=alphabetic,
]{biblatex}
\addbibresource{bibliography.bib} %Imports bibliography file

%%%%%%%%%%
% Title  %
%%%%%%%%%%
\title{
אלגברה ב' - גיליון תרגילי בית 2 \\
הדטרמיננטה
\\
\vspace{1cm}
\large{תאריך הגשה: 30.4.2026}
}
\date{}

\begin{document}
\maketitle

\begin{exercise}
חשבו את הדטרמיננטות של המטריצות הבאות.

\begin{enumerate}
\item
\begin{align*}
A &\coloneqq \pmat{
0 & 0 & 1 & 2 \\
0 & 0 & 3 & 4 \\
1 & 2 & 0 & 0 \\
3 & 4 & 0 & 0
} \in \Mat_4\prs{\RR}
\\
B &\coloneqq \pmat{
1 & 1 & 2 & 2 \\
1 & 0 & 2 & 0 \\
3 & 3 & 0 & 0 \\
3 & 0 & 0 & 0
} \in \Mat_4\prs{\RR}
\\
C &\coloneqq \pmat{
0 & \cdots & \cdots & 0 & 1 \\
1 & 0 & \ddots & \ddots & 0 \\
0 & \ddots & \ddots & \ddots & \vdots \\
\vdots & \ddots & \ddots & 0 & \vdots \\
0 & \cdots & 0 & 1 & 0
} \in \Mat_n\prs{\RR}
\\
D &\coloneqq \pmat{
1 & 1 & 1 \\
1 & i & i^2 \\
1 & -i & \prs{-i}^2
} \in \Mat_3\prs{\CC}
\end{align*}
\end{enumerate}
\end{exercise}

\begin{exercise}
חשבו את הדטרמיננטות של האופרטורים הבאים.
\begin{enumerate}
\item
\begin{align*}
T \colon \Mat_2\prs{\RR} &\to \Mat_2\prs{\RR} \\
A &\mapsto A^t
\end{align*}
\item
\begin{align*}
T \colon \Mat_2\prs{\RR} &\to \Mat_2\prs{\RR} \\
A &\mapsto \pmat{0 & 1 \\ 1 & 0} A
\end{align*}
\item
\begin{align*}
T \colon \CC_2\brs{x} &\to \CC_2\brs{x} \\
p\prs{x} &\mapsto p'(x)
\end{align*}
\end{enumerate}
\end{exercise}

\begin{exercise}
\begin{enumerate}
\item
תהי
\begin{equation}\label{eq:block-diag}
A = \pmat{
A_1 & 0 & \cdots & 0 & 0 \\
0 & A_2 & \ddots & \ddots & 0 \\
\vdots & \ddots & \ddots & \ddots & \vdots \\
0 & \ddots & \ddots & A_{n-1} & 0 \\
0 & 0 & \cdots & 0 & A_k
} \in \Mat_n\prs{\FF}
\end{equation}
עבור
$A_i \in \Mat_{m_i}\prs{\FF}$
ועבור
$m_i \in \NN_+$.
נזכיר, מהתרגול, כי
\begin{align*}
\text{.} \det\prs{A} = \prod_{i \in [k]} \det\prs{A_i}
\end{align*}
חשבו את הדטרמיננטות של המטריצות הבאות.
\begin{enumerate}
\item
\[A_i\prs{\alpha} \coloneqq \pmat{
A_1 & 0 & \cdots & \cdots & \cdots & \cdots & \cdots & 0 & 0 \\
0 & A_2 & \ddots & \ddots & \ddots & \ddots & \ddots & \ddots &  0 \\
\vdots & \ddots & \ddots & \ddots & \ddots & \ddots & \ddots & \ddots & \vdots \\
\vdots & \ddots & \ddots & A_{i-1} & \ddots & \ddots & \ddots & \ddots & \vdots \\
\vdots & \ddots & \ddots & \ddots & \alpha A_i & \ddots & \ddots & \ddots & \vdots \\
\vdots & \ddots & \ddots & \ddots & \ddots & A_{i+1} & \ddots & \ddots & \vdots \\
\vdots & \ddots & \ddots & \ddots & \ddots & \ddots & \ddots & \ddots & \vdots \\
0 & \ddots & \ddots &\ddots & \ddots & \ddots & \ddots & A_{n-1} & 0 \\
0 & 0 & \cdots & \cdots & \cdots & \cdots & \cdots & 0 & A_k
}
\in \Mat_n\prs{\FF}\]
עבור
$\alpha \in \FF$.
\item
\[A_{i \leftrightarrow j} \coloneqq \pmat{
A_1 & 0 & \cdots & \cdots & \cdots & \cdots & \cdots & \cdots & 0 & 0 \\
0 & A_2 & \ddots & \ddots & \ddots & \ddots & \ddots & \ddots & \ddots &  0 \\
\vdots & \ddots & \ddots & \ddots & \ddots & \ddots & \ddots & \ddots & \ddots & \vdots \\
\vdots & \ddots & \ddots & A_{i-1} & \ddots & \ddots & \ddots & \ddots & \vdots \\
\vdots & \ddots & \ddots & \ddots & 0 & A_{i+1} & \ddots & \ddots & \vdots \\
\vdots & \ddots & \ddots & \ddots & A_{i} & 0 & \ddots & \ddots & \vdots \\
\vdots & \ddots & \ddots & \ddots & \ddots & \ddots & A_{i+2} & \ddots & \ddots & \vdots \\
\vdots & \ddots & \ddots & \ddots & \ddots & \ddots & \ddots & \ddots & \ddots & \vdots \\
0 & \ddots & \ddots &\ddots & \ddots & \ddots & \ddots & \ddots & A_{n-1} & 0 \\
0 & 0 & \cdots & \cdots & \cdots & \cdots & \cdots & \cdots & 0 & A_k
}
\in \Mat_n\prs{\FF}\]
המתקבלת על ידי החלפת השורות ה־%
$i,i+1$
ב־%
\eqref{eq:block-diag}.

\textbf{רמז:}
חשבו את הדטרמיננטה של המטריצה
$\pmat{0 & A_{i+1} \\
A_i & 0}$.

\item
\[A_{i,j}\prs{\alpha} \coloneqq \pmat{
A_1 & 0 & \cdots & \cdots & \cdots & \cdots & \cdots & 0 & 0 \\
0 & A_2 & \ddots & \ddots & \ddots & \ddots & \ddots & \ddots &  0 \\
\vdots & \ddots & \ddots & \ddots & \ddots & \ddots & \ddots & \ddots & \vdots \\
\vdots & \ddots & \ddots & A_{i} & \ddots & \ddots & \ddots & \ddots & \vdots \\
\vdots & \ddots & \ddots & \ddots & \ddots & \ddots & \ddots & \ddots & \vdots \\
\vdots & \ddots & \ddots & \alpha A_i & \ddots & A_{j} & \ddots & \ddots & \vdots \\
\vdots & \ddots & \ddots & \ddots & \ddots & \ddots & \ddots & \ddots & \vdots \\
0 & \ddots & \ddots &\ddots & \ddots & \ddots & \ddots & A_{n-1} & 0 \\
0 & 0 & \cdots & \cdots & \cdots & \cdots & \cdots & 0 & A_k
}
\in \Mat_n\prs{\FF}\]
המתקבלת על ידי הוספת
$\alpha$
כפול השורה ה־%
$i$
ב־%
\eqref{eq:block-diag}
לשורה ה־%
$j$.
\end{enumerate}
\item
היעזרו בסעיף הקודם כדי לחשב את הדטרמיננטות של המטריצות הבאות מהתרגיל הראשון.
\begin{align*}
A &= \pmat{
0 & 0 & 1 & 2 \\
0 & 0 & 3 & 4 \\
1 & 2 & 0 & 0 \\
3 & 4 & 0 & 0
} \in \Mat_4\prs{\RR}
\\
B &= \pmat{
1 & 1 & 2 & 2 \\
1 & 0 & 2 & 0 \\
3 & 3 & 0 & 0 \\
3 & 0 & 0 & 0
} \in \Mat_4\prs{\RR}
\\
C &= \pmat{
0 & \cdots & \cdots & 0 & 1 \\
1 & 0 & \ddots & \ddots & 0 \\
0 & \ddots & \ddots & \ddots & \vdots \\
\vdots & \ddots & \ddots & 0 & \vdots \\
0 & \cdots & 0 & 1 & 0
} \in \Mat_n\prs{\RR}
\end{align*}
\end{enumerate}
\end{exercise}

\end{document}